\documentclass[11pt]{article}
\usepackage[letterpaper,margin=0.78in]{geometry}
\usepackage{amsmath,amssymb,mathtools}
\usepackage{array,booktabs}
\usepackage{microtype}
\usepackage{enumitem}
\usepackage[T1]{fontenc}
\usepackage{lmodern}
\usepackage{fancyhdr}
\usepackage{xcolor}
\usepackage{tcolorbox}
\usepackage{hyperref}
\hypersetup{hidelinks}
\setlength{\parindent}{0pt}
\setlength{\parskip}{6pt}
\setlist[itemize]{leftmargin=1.4em,itemsep=2pt,topsep=2pt}
\pagestyle{fancy}
\fancyhf{}
\lhead{LANGUAGE MECHANICS}
\rhead{IDENTITY \& REFERENCE v0.2}
\cfoot{\thepage}
\setlength{\headheight}{14pt}
\newtcolorbox{lawbox}{colback=black!3,colframe=black!65,boxrule=0.45pt,arc=0pt,left=8pt,right=8pt,top=7pt,bottom=7pt}
\newcommand{\Dom}{\operatorname{Dom}}
\newcommand{\ID}{\mathsf{ID}}

\begin{document}

\begin{center}
{\Large\bfseries IDENTITY \& REFERENCE}\\[3pt]
{\large Reapplication, Quotient Identity, and Active Reference}\\[8pt]
Anthony Vito Coppa\\[2pt]
\textit{Original construction sequence: 23--26 January 2026}\\
\textit{Canonical standalone edition v0.2: 29 August 2026}
\end{center}

\begin{lawbox}
\textbf{Abstract.}
This paper separates three formal questions that ordinary language often compresses into sameness. Reapplication asks whether a partial operation remains defined when used again. Identity asks whether the result remains in the same declared reader class. Reference selects which already-admissible reader class is held fixed when more than one class remains available. A reader-induced equivalence relation makes the quotient type explicit. When a partial operation respects that relation on its admitted domain, it descends to a map on quotient classes, and operational identity is exactly a quotient fixed point. Exact state equality is stronger. Reference is therefore neither identity nor a source of identity: it is a versioned selection among already-admissible identity classes. The construction requires no substance, time, semantics, physical realization, or theory of meaning.
\end{lawbox}

\section{Scope}

Fix a candidate class $\Sigma$ and a declared reader
\[
\rho:\Sigma\to Y.
\]
The reader induces the equivalence relation
\[
\boxed{x\sim_\rho y\iff \rho(x)=\rho(y).}
\]
Reflexivity, symmetry, and transitivity follow from equality in $Y$, so the quotient $\Sigma/{\sim_\rho}$ is well-defined.

Let
\[
T:\Dom(T)\subseteq\Sigma\longrightarrow\Sigma
\]
be a partial operation admitted for study. No identity claim follows merely from the existence of $T$.

The definitions in this paper are \textbf{DECLARED}. The theorem, corollary, and proposition below are \textbf{DERIVED} from those declarations. No additional document-level status is imposed.

\begin{lawbox}
\centering
\textbf{Question.} When the same formal operation may be used again, what must be distinguished in order to state identity and reference exactly?
\end{lawbox}

\section{Reapplication}

\textbf{Definition 2.1 (Reapplication).}
For $x\in\Dom(T)$, a second application is admissible at $x$ exactly when
\[
T(x)\in\Dom(T).
\]
Equivalently, $T^2(x)$ is defined.

Reapplication says only \emph{again}. It supplies no duration, history, explanation, or guarantee of indefinite continuation.

\textbf{Definition 2.2 (Finite and indefinite reuse).}
For $n\ge1$, reuse through $n$ is admissible at $x$ when $T^k(x)$ is defined for every $1\le k\le n$. Reuse is indefinitely admissible at $x$ when $T^n(x)$ is defined for every finite $n$.

\section{Recognition-stable identity}

Reapplication and identity are different tests. Reapplication asks whether the operation remains defined. Identity asks whether the declared reader class survives the lawful move.

\textbf{Definition 3.1 (Operational identity under reuse).}
For $x\in\Dom(T)$, identity under one admitted application holds when
\[
\boxed{T(x)\sim_\rho x.}
\]
Identity under reuse through $n$ holds when all required iterates are defined and
\[
T^k(x)\sim_\rho x
\qquad(1\le k\le n).
\]

This is reader-relative identity. Exact state return
\[
T(x)=x
\]
is stronger and is never substituted silently for $T(x)\sim_\rho x$.

\section{The quotient fixed point}

Restrict $\sim_\rho$ to $\Dom(T)$. Assume $T$ respects the reader equivalence on its admitted domain:
\[
x\sim_\rho y
\quad\Longrightarrow\quad
T(x)\sim_\rho T(y)
\qquad(x,y\in\Dom(T)).
\]
Then the map
\[
\overline T:\Dom(T)/{\sim_\rho}\longrightarrow \Sigma/{\sim_\rho},
\qquad
\overline T([x]_\rho)=[T(x)]_\rho
\]
is well-defined.

\textbf{Theorem 4.1 (Reuse fixed point).}
For every $x\in\Dom(T)$,
\[
\boxed{
T(x)\sim_\rho x
\iff
\overline T([x]_\rho)=[x]_\rho.
}
\]
Thus operational identity under reuse is exactly a fixed-point condition at the declared quotient reader.

\textit{Proof.}
By definition,
\[
\overline T([x]_\rho)=[T(x)]_\rho.
\]
Equality of quotient classes satisfies
\[
[T(x)]_\rho=[x]_\rho
\iff
T(x)\sim_\rho x.
\]
The two statements are therefore equivalent. $\square$

\textbf{Corollary 4.2 (Exact fixed point is stronger).}
If $T(x)=x$, then $T(x)\sim_\rho x$. Conversely, reader identity implies exact equality only when the reader is injective on the admitted comparison set containing $x$ and $T(x)$.

\section{Closure and idempotence}

A related construction occurs when an operation is intended to stabilize a candidate under repeated use.

\textbf{Definition 5.1 (Recognition-idempotent closure).}
Let $C$ be a partial map with $x,C(x)\in\Dom(C)$. It is recognition-idempotent at $x$ when
\[
\boxed{C(C(x))\sim_\rho C(x).}
\]
Exact idempotence
\[
C(C(x))=C(x)
\]
is stronger.

Recognition-idempotence states stabilization under the declared reader. It does not define substance, selfhood, memory, or meaning.

\section{Multiplicity and active reference}

Reader identity need not determine a unique baseline. Let
\[
\ID\subseteq \Sigma/{\sim_\rho}
\]
be the set of reader classes presently admitted for a stated comparison.

\textbf{Definition 6.1 (Active reference).}
An active reference is a declared selection
\[
r\in\ID
\]
held fixed for a stated comparison or version. A candidate $x$ returns to that reference exactly when
\[
[x]_\rho=r.
\]

Reference performs a selection office. It fixes which already-admissible identity class is in force for the stated use. It does not create that class and does not establish why the class should be preferred outside the declared use.

\textbf{Proposition 6.2 (Reference requirement under multiplicity).}
If $|\ID|>1$, the phrase ``the identity in force'' is underdetermined until an active reference $r\in\ID$ is declared.

\textit{Proof.}
More than one admitted reader class satisfies the upstream requirements. Without a selection, no unique class is designated as the comparison baseline. Declaring $r$ supplies that missing address and nothing further. $\square$

Changing the active reference creates a new reference version. It does not retroactively alter results obtained under an earlier reference.

\section{Identity and reference are different offices}

\begin{center}
\begin{tabular}{@{}lll@{}}
\toprule
\textbf{Pressure} & \textbf{Formal operation} & \textbf{Office}\\
\midrule
Reuse & $T(x)\sim_\rho x$ & reader-relative identity\\
Multiplicity & choose $r\in\ID$ & active reference\\
\bottomrule
\end{tabular}
\end{center}

The same ordinary word \emph{sameness} can appear in both pressures. The formal offices remain separate:
\[
\boxed{
\text{identity tests preservation under reuse; reference fixes the baseline under alternatives.}
}
\]

\section{Boundary}

This paper establishes only:
\begin{itemize}
\item admissible reapplication for a partial operation;
\item reader-relative identity under reuse;
\item the quotient fixed-point characterization when the operation respects reader equivalence;
\item recognition-idempotent closure;
\item reference as declared fixation of an already-admissible identity class;
\item versioning when the active reference changes.
\end{itemize}

It does not establish why a reference should be selected, a theory of truth or meaning, a physical or biological realization, memory as a receipt-and-decoder office, or a claim that one local reuse grammar governs every jurisdiction.

\section{Canonical statement}

\begin{lawbox}
Reapplication asks whether the operation remains defined. Identity asks whether the declared reader class survives that lawful use. When the operation respects the reader equivalence, identity is exactly a fixed point of the induced quotient map. If more than one reader class remains admitted for a comparison, reference is the declared fixation of the class in force. Identity preserves under reuse; reference selects under multiplicity.
\end{lawbox}

\section*{Provenance}

The construction consolidated here was developed across three January records: \emph{Identity \& Identity: The Rule of Reference} (23 January 2026), \emph{Application \& Reuse: Identity Fixed-point} (24 January 2026), and \emph{Identity \& Reference: Reuse Fixed Point} (26 January 2026). These titles and dates are retained as provenance only. The present standalone form contains every definition and proof used by its claims.

\end{document}
